% ==============================================================================
% T0-Theorie: Dokument 167
% Titel: LiNbO₃-Materialeigenschaften auf der ξ-Potenzleiter —
%        Phasenanpassung in PPLN: Prüfung einer geometrischen Zuordnung
% Autor: Johann Pascher
% Datum: März 2026
% Bezug: Dokumente 041, 056, 061, 081, 122, 133, 161, 162, 165, 166
% ==============================================================================

\section*{Abstract}

Periodisch gepoltes Lithiumniobat (PPLN) ist der physikalische Kern der
Kohärenten Ising-Maschine (CIM): Der $\chi^2$-Nichtlinearitätskristall
implementiert durch parametrische Verstärkung den DOPO-Spin, das fundamentale
Bit des Ising-Rechners. Die Polungsperiode $\Lambda$, der Brechungsindex
$n(\omega)$ und der elektrooptische Koeffizient $r_{33}$ gelten als empirisch
bestimmte Materialparameter ohne geometrischen Ursprung.

Dieses Dokument prüft, ob ein zentraler Parameter von LiNbO$_3$, der
Dispersionsunterschied im DOPO-Betrieb, auf der $\xipar$-Potenzleiter der
T0-Theorie liegt, und zwar in der Form:

\begin{equation}
    \Delta n_{(1064 \to 2128\,\text{nm})} \;\stackrel{?}{=}\; \sqrt{3}\cdot\xipar^{1/2} = 0{,}020
    \label{eq:main}
\end{equation}

mit $\xipar = 4/30\,000$. Diese Zuordnung gilt nicht: Die Sellmeier-Gleichung nach
Zelmon 1997 (kongruentes LiNbO$_3$, außerordentlich) gibt
$n_e(1064\,\text{nm}) = 2{,}1555$ und $n_e(2128\,\text{nm}) = 2{,}1219$, also
$\Delta n = 0{,}0336$. Nur dieser Wert passt zur Polungsperiode:
$\lambda_p/\Delta n = 1{,}064\,\mu\text{m}/0{,}0336 = 31{,}6\,\mu$m ($\Lambda = 31\,\mu$m),
während $0{,}020$ $\Lambda = 53\,\mu$m verlangen würde. $\sqrt{3}\,\xipar^{1/2} = 0{,}020$
liegt damit $41\,\%$ unter dem Messwert. Die exakte Gleichheit
$\sqrt{3}\cdot(4/30000)^{1/2} = \sqrt{12/30000} = 0{,}02$ ist eine algebraische
Identität; der Messwert ist nur auf 2--3 Stellen bestimmt.

Der Exponent $1/2$ erscheint in der T0-Theorie als schwache
Kopplungskonstante $\alpha_W = \xipar^{1/2}$, als Leptonmassen-Verhältnis
$m_e/m_\mu \propto \xipar^{1/2}$ und an weiteren Stellen
(Dok.~041, 056, 061, 122, 133). Die Relationen $E_\mathrm{EW} \sim E_P \cdot \xipar^{1/2}$
und $m_h/M_P = \xipar^{1/2}$ sind nicht tragfähig. Die PPLN-Dispersion liegt nicht
auf dieser Stufe; als Treffer auf der $\xipar$-Leiter bleibt das Verhältnis der
elektrooptischen Koeffizienten $r_{33}/d_{33} \approx 1 + \xipar^{1/4}$
(Abweichung $2{,}2\,\%$).

\clearpage

% ==============================================================================
\section{Einleitung: PPLN als Kern der Kohärenten Ising-Maschine}
\label{sec:intro}
% ==============================================================================

\subsection{Der DOPO-Mechanismus}

Die Kohärente Ising-Maschine (CIM) implementiert das Ising-Hamiltonproblem
$H = -\sum_{ij} J_{ij}\sigma_i\sigma_j$ durch ein Netzwerk von DOPO-Pulsen
(Degenerate Optical Parametric Oscillators) in einer Faserkavität.

Der physikalische Kern ist der $\chi^2$-Nichtlinearitätskristall: Ein Pumpphoton
der Frequenz $\omega_p$ wird in zwei Signal- und Idlerphotonen bei $\omega_p/2$
umgewandelt. Oberhalb der Schwelle entwickelt der DOPO exakt zwei stabile
Zustände — Phase 0 oder $\pi$ — die dem Ising-Spin $\sigma_i = \pm 1$
entsprechen.

Das Material der Wahl ist periodisch gepoltes Lithiumniobat (PPLN). Die Frage,
warum LiNbO$_3$, wird üblicherweise rein empirisch beantwortet: stärkster
bekannter $d_{33}$-Koeffizient bei niedrigen Absorptionsverlusten.
Dieses Dokument prüft, ob sich dafür eine geometrische Antwort aus der
$\xipar$-Potenzleiter ergibt.

\subsection{Die Phasenanpassungsbedingung}

Im quasi-Phasenanpassungsschema (QPM) wird die Phasenfehlanpassung
$\Delta k = k_p - k_s - k_i \neq 0$ durch periodisches Umpolen der
Kristalldomänen mit Periode $\Lambda$ kompensiert:

\begin{equation}
    \Lambda = \frac{2\pi}{\Delta k} = \frac{\lambda_{\text{pump}}}{\Delta n}
    \label{eq:poling}
\end{equation}

Die Polungsperiode ist direkt proportional zum inversen
Dispersionsunterschied $\Delta n = n(\lambda_{\text{pump}}) -
n(2\lambda_{\text{pump}})$ ($\lambda_{\text{pump}}$: Pumpwellenlänge beim DOPO, SH-Wellenlänge bei SHG).
Für den DOPO ist $\Delta k = 2\pi(n_p - n_s)/\lambda_p$, also $\Lambda = \lambda_p/(n_p - n_s)$.

\begin{table}[H]
    \centering
    \adjustbox{max width=\textwidth}{%
\begin{tabular}{lccc}
        \toprule
        \textbf{Prozess} & \textbf{Pumpe} &
        \textbf{$\Lambda$\,($\mu$m)} & \textbf{$\Delta n$} \\
        \midrule
        OPO — DOPO-Betrieb (CIM)  & $1064\,\text{nm}$ & $31{,}0$ & $0{,}034$ \\
        SHG — Telekommunikation   & $780\,\text{nm}$  & $19{,}5$ & $0{,}040$ \\
        SHG — Grüner Laser        & $532\,\text{nm}$  & $6{,}8$ & $0{,}078$ \\
        \bottomrule
    \end{tabular}%
}
    \caption{Typische PPLN-Parameter für LiNbO$_3$ (außerordentliche
    Polarisation, Raumtemperatur, Sellmeier nach Zelmon 1997).}
    \label{tab:ppln}
\end{table}
Nach Sellmeier (Zelmon 1997, $n_e$) gelten im Einzelnen: DOPO $1064 \to 2128$~nm
$\Delta n = 0{,}0336$, $\Lambda = 31{,}6\,\mu$m; SHG $1560 \to 780$~nm $\Delta n = 0{,}0405$,
$\Lambda = 19{,}3\,\mu$m; SHG $1064 \to 532$~nm $\Delta n = 0{,}078$,
$\Lambda = 0{,}532/0{,}078 = 6{,}8\,\mu$m.

\clearpage

% ==============================================================================
\section{%
  Der Exponent \texorpdfstring{$\xipar^{1/2}$}{xi\^{}(1/2)} in der T0-Theorie}
\label{sec:xi_halb}
% ==============================================================================

\subsection{%
  \texorpdfstring{$\xipar^{1/2}$}{xi\^{}(1/2)} als fundamentaler
  elektroschwacher Exponent}

In der T0-Theorie ist $\xipar^{1/2} = (4/30\,000)^{1/2} = 0{,}011547\ldots$
nicht eine unter vielen $\xipar$-Potenzen, sondern der fundamentale
Übergangsexponent zwischen Planck-Skala und Teilchenphysik. Er tritt in
einer bemerkenswerten Vielzahl unabhängiger Kontexte auf:

\begin{table}[H]
    \centering
    \adjustbox{max width=\textwidth}{%
\begin{tabular}{llll}
        \toprule
        \textbf{Physikalische Größe} & \textbf{T0-Ausdruck} &
        \textbf{Wert} & \textbf{Dok.} \\
        \midrule
        Schwache Kopplungskonstante &
            $\alpha_W = \xipar^{1/2}$ &
            $1{,}15 \times 10^{-2}$ & 041, 061 \\
        Leptonmassen-Verhältnis &
            $m_e/m_\mu \propto \xipar^{1/2}$ &
            $4{,}84 \times 10^{-3}$ & 056, 122 \\
        Bottom-Quark-Masse &
            $m_b = v \cdot \tfrac{3}{2}\cdot\xipar^{1/2}$ &
            $4{,}25\,\text{GeV}$ & 041 \\
        Planck-Konstante (Skalierung) &
            $\hbar \sim \sqrt{\xipar}$ & — & 105 \\
        \bottomrule
    \end{tabular}%
}
    \caption{Der Exponent $\xipar^{1/2} = 0{,}011547$ tritt in der T0-Theorie
    als universeller elektroschwacher Exponent auf. Er markiert den geometrischen
    Übergang von der Planck-Welt zur Teilchenphysik-Skala.}
    \label{tab:xi12}
\end{table}

Zwei weitere $\xipar^{1/2}$-Relationen gelten nicht: $m_h/M_P = \xipar^{1/2}$ ergibt
$1{,}15\times10^{-2}$, gemessen ist $m_h/M_P = 1{,}0\times10^{-17}$ (siehe Dok.~068);
$E_\mathrm{EW} \sim E_P\cdot\xipar^{1/2}$ ergibt mit $E_P = 1{,}22\times10^{19}$~GeV etwa
$1{,}4\times10^{17}$~GeV, nicht $\sim 100$~GeV, und ist damit als Herleitung der
elektroschwachen Skala nicht tragfähig.

\subsection{Physikalische Bedeutung des Exponenten}

Der Exponent $1/2$ markiert in der T0-Geometrie den \emph{elektroschwachen
Übergang} auf der $\xipar$-Potenzleiter:

\begin{align}
    \text{Planck-Skala:}\quad          & \xipar^0     = 1
        && E \sim E_P \notag \\
    \text{Elektroschwacher Übergang:}\quad & \xipar^{1/2} \approx 10^{-2}
        && \alpha_W,\; m_e/m_\mu
        \label{eq:stufe} \\
    \text{Teilchenskala:}\quad         & \xipar^1     \approx 10^{-4}
        && m_e,\; \alpha_\mathrm{EM} \notag
\end{align}

Die PPLN-Dispersion liegt nicht auf dieser Stufe (Abschnitt~\ref{sec:hauptergebnis}).

\subsection{Abgrenzung zu Dokument 166}

Dokument~166 verwendet für die Casimir-Skala den Exponenten $1/4$
($L_\xi \propto \xipar^{1/4}$, Zerlegung $41/4 = 10 + 1/4$), nicht $\xipar^{1/2}$;
einen gemeinsamen Exponenten von Casimir-Skala und PPLN-Dispersion gibt es nicht.

\clearpage

% ==============================================================================
\section{Prüfung: \texorpdfstring{$\Delta n$ gegen $\sqrt{3}\cdot\xipar^{1/2}$}{%
  Delta n vs. sqrt(3) xi\^{}(1/2)}}
\label{sec:hauptergebnis}
% ==============================================================================

\subsection{Numerischer Vergleich}

Der Dispersionsunterschied von LiNbO$_3$ für den DOPO-Betrieb berechnet
sich aus den Sellmeier-Brechungsindizes der außerordentlichen Polarisation
(Zelmon 1997):

\begin{equation}
    \Delta n = n_e(1064\,\text{nm}) - n_e(2128\,\text{nm})
    = 2{,}1555 - 2{,}1219 = 0{,}0336
\end{equation}

Der Ausdruck aus Gl.~\eqref{eq:main} ergibt:

\begin{align}
    \sqrt{3}\cdot\xipar^{1/2}
    &= \sqrt{12/30000} = 0{,}020000\ldots
\end{align}

\begin{center}
\adjustbox{max width=\textwidth}{%
\begin{tabular}{lc}
    \toprule
    & \textbf{Wert} \\
    \midrule
    Gemessen: $\Delta n(1064 \to 2128\,\text{nm})$ & $0{,}0336$ \\
    T0-Ausdruck: $\sqrt{3}\cdot\xipar^{1/2}$       & $0{,}0200$ \\
    \bottomrule
\end{tabular}%
}
\end{center}

\medskip
Die Zuordnung $\Delta n = \sqrt{3}\cdot\xipar^{1/2}$ gilt nicht, weil der Ausdruck
$41\,\%$ unter dem Messwert liegt. Die exakte Gleichheit mit $0{,}020$ ist eine
algebraische Identität; der Messwert ist nur auf 2--3 Stellen bestimmt.

\clearpage

% ==============================================================================
\section{Weitere Treffer auf der \texorpdfstring{$\xipar$}{xi}-Leiter}
\label{sec:weitere}
% ==============================================================================

\subsection{Verhältnis der elektrooptischen Koeffizienten}

LiNbO$_3$ besitzt zwei relevante Kopplungskonstanten: den linearen
elektrooptischen (Pockels-)Koeffizienten $r_{33} = 30{,}8\,\text{pm/V}$
und den nichtlinearen Koeffizienten $d_{33} = 27{,}2\,\text{pm/V}$.
Ihr Verhältnis:

\begin{equation}
    \frac{r_{33}}{d_{33}} = \frac{30{,}8}{27{,}2} = 1{,}1324
    \;\approx\; 1 + \xipar^{1/4} = 1{,}1075
    \qquad (\text{Abw.}\ {+}2{,}2\,\%)
\end{equation}

$\xipar^{1/4}$ ist in T0 der Sub-Planck-Geometrieexponent (Dok.~009).
Das Verhältnis zweier elektrooptischer Kopplungskonstanten liegt auf der
$\xipar$-Leiter — ein Hinweis, dass auch die Kristallbandstruktur selbst
$\xipar$-rational ist.

\clearpage

% ==============================================================================
\section{Einordnung: \texorpdfstring{$\xipar^{1/2}$}{xi\^{}(1/2)} als
  elektroschwache Stufe}
\label{sec:einordnung}
% ==============================================================================

Die folgende Tabelle zeigt die vollständige Einordnung von $\xipar^{1/2}$
über alle Bereiche der T0-Theorie:

\begin{table}[H]
    \centering
    \adjustbox{max width=\textwidth}{%
\begin{tabular}{llll}
        \toprule
        \textbf{Bereich} & \textbf{Größe} &
        \textbf{T0-Ausdruck} & \textbf{Dok.} \\
        \midrule
        \multicolumn{4}{l}{\itshape Teilchenphysik (Standardmodell)} \\[2pt]
        & Schwache Kopplung   & $\alpha_W = \xipar^{1/2}$
            & 041, 061 \\
        & Leptonmassen        & $m_e/m_\mu \propto \xipar^{1/2}$
            & 056, 122, 133 \\
        & Bottom-Quark        & $m_b = v\cdot\tfrac{3}{2}\cdot\xipar^{1/2}$
            & 041 \\
        \midrule
        \multicolumn{4}{l}{\itshape Materialeigenschaften (dieses Dokument)} \\[2pt]
        & EO-Koeffizienten
            & $r_{33}/d_{33} \approx 1 + \xipar^{1/4}$
            & \textbf{167} \\
        \midrule
        \multicolumn{4}{l}{\itshape Kosmologie und Quantenoptik} \\[2pt]
        & Casimir-Skala       & $L_\xi = (15\xipar/\pi^2)^{1/4}\,\hbar c/(k_BT_\text{CMB})$
            & 078, 166 \\
        & $H_0$-Verhältnis    & $\hbar H_0/m_ec^2 = (\pi/2)\cdot\xipar^{10}$
            & 165 \\
        \bottomrule
    \end{tabular}%
}
    \caption{$\xipar^{1/2}$ als elektroschwacher Exponent und weitere
    $\xipar$-Potenzen der T0-Theorie. Die PPLN-Dispersion liegt nicht auf der
    $\xipar^{1/2}$-Stufe (Abschnitt~\ref{sec:hauptergebnis}).}
    \label{tab:universal}
\end{table}

Die PPLN-Dispersion gehört nicht in diese Hierarchie: $\sqrt{3}\cdot\xipar^{1/2}$
liegt $41\,\%$ unter dem gemessenen $\Delta n = 0{,}0336$. Ebenso fehlen die
Relationen $E_\mathrm{EW} \sim E_P\cdot\xipar^{1/2}$ und $m_h/M_P = \xipar^{1/2}$, weil
sie nicht tragfähig sind (Abschnitt~\ref{sec:xi_halb}).

\clearpage

% ==============================================================================
\section{Konsequenzen für das CIM-Design}
\label{sec:konsequenzen}
% ==============================================================================

Eine geometrische Antwort auf „Warum LiNbO$_3$?" über
$\Delta n = \sqrt{3}\cdot\xipar^{1/2}$ gilt nicht, weil die gemessene Dispersion
$0{,}0336$ beträgt. Die Materialwahl bleibt empirisch begründet: stärkster
bekannter $d_{33}$-Koeffizient bei niedrigen Absorptionsverlusten. Eine auf
diese Zuordnung gestützte Auswahlregel $\Delta n = k\cdot\xipar^{n/4}$ für
CIM-Substrate und Vorhersagen daraus ($\sqrt{2}\cdot\xipar^{1/2}$ bei
$\sim 900\,\text{nm}$, $2\cdot\xipar^{1/2}$ bei $\sim 680\,\text{nm}$) gelten daher nicht.

\clearpage

% ==============================================================================
\section{Zusammenfassung}
\label{sec:zusammenfassung}
% ==============================================================================

\begin{enumerate}

    \item \textbf{Ergebnis}: Die Dispersion von LiNbO$_3$ für den
    wichtigsten PPLN-Prozess ($1064 \to 2128\,\text{nm}$) beträgt nach
    Sellmeier (Zelmon 1997) $\Delta n = 0{,}0336$. Die Zuordnung
    $\Delta n = \sqrt{3}\cdot\xipar^{1/2} = 0{,}020$ gilt nicht: Der Ausdruck liegt
    $41\,\%$ unter dem Messwert, die Gleichheit mit $0{,}020$ ist eine Identität.

    \item \textbf{Exponent}: $\xipar^{1/2}$ erscheint als
    $\alpha_W = \xipar^{1/2}$, $m_e/m_\mu \propto \xipar^{1/2}$ und in weiteren
    Größen (Dok.~041, 056, 061, 122, 133). $E_\mathrm{EW} \sim E_P\cdot\xipar^{1/2}$
    und $m_h/M_P = \xipar^{1/2}$ sind nicht tragfähig. Mit der Casimir-Skala, die
    in Dok.~166 den Exponenten $1/4$ trägt, teilt $\xipar^{1/2}$ keinen Exponenten.

    \item \textbf{Weiterer Treffer}: $r_{33}/d_{33} \approx 1 + \xipar^{1/4}$
    (Abw.\ $2{,}2\,\%$).

    \item \textbf{Einordnung und Schlussfolgerung}: Die PPLN-Dispersion liegt
    nicht auf der $\xipar^{1/2}$-Stufe. Die Materialwahl für das CIM-Design bleibt
    empirisch begründet.

\end{enumerate}

\bigskip\bigskip

\noindent\textbf{Verwandte Dokumente:}
Dok.~041 (Parameterherleitung, $\alpha_W = \xipar^{1/2}$),
Dok.~056 (Universale Ableitung, Leptonmassen),
Dok.~061 (CMB-Einheiten, $\alpha_W$),
Dok.~081 (Zusammenfassung, elektroschwache Skala),
Dok.~122 (Verhältnis–absolut, $m_e/m_\mu$),
Dok.~133 (Fraktale Korrektur),
Dok.~161 (CIM-Grundlagen),
Dok.~165 ($H_0$-Verhältnis),
Dok.~166 (Casimir-CMB-$H_0$-Brücke).
